% =========================================================
% ppt/src/minggu-03.tex
% Slide pertemuan 3
% Bab V — Register 8086 dan Register Flag
%
% Sub-CPMK 1.2 | CPMK01 | CPL02
% =========================================================

\documentclass[aspectratio=169,11pt]{beamer}
% =========================================================
% ppt/src/preamble-ppt.tex
% Tema, warna, paket, dan makro bersama seluruh deck slide.
%
% Dipakai oleh setiap berkas minggu-NN.tex dengan pola:
%   \documentclass[aspectratio=169,11pt]{beamer}
%   % =========================================================
% ppt/src/preamble-ppt.tex
% Tema, warna, paket, dan makro bersama seluruh deck slide.
%
% Dipakai oleh setiap berkas minggu-NN.tex dengan pola:
%   \documentclass[aspectratio=169,11pt]{beamer}
%   % =========================================================
% ppt/src/preamble-ppt.tex
% Tema, warna, paket, dan makro bersama seluruh deck slide.
%
% Dipakai oleh setiap berkas minggu-NN.tex dengan pola:
%   \documentclass[aspectratio=169,11pt]{beamer}
%   \input{preamble-ppt}
%
% Palet warna disalin dari buku_ajar/config/styles.tex agar
% identitas visual slide sama dengan buku ajar.
% =========================================================

% ---------------------------------------------------------
% 1. Paket
% ---------------------------------------------------------
\usepackage[T1]{fontenc}
\usepackage[utf8]{inputenc}
\usepackage[indonesian]{babel}
\usepackage{lmodern}
\usepackage{tikz}
\usetikzlibrary{arrows.meta, positioning, calc, shapes.geometric}
\usepackage{booktabs}
\usepackage{array}
\usepackage{amsmath}
\usepackage{amsfonts}
\usepackage{pifont}          % \ding untuk tanda checklist
\usepackage{listings}

% ---------------------------------------------------------
% 2. Palet warna OBE (identik dengan buku)
% ---------------------------------------------------------
\definecolor{obeBlue}{HTML}{1F3864}
\definecolor{obeBlueLight}{HTML}{DCE4F5}
\definecolor{obeTeal}{HTML}{0F6E6E}
\definecolor{obeTealLight}{HTML}{D9F0EE}
\definecolor{obeGreen}{HTML}{2E6B34}
\definecolor{obeGreenLight}{HTML}{DEEFDF}
\definecolor{obeAmber}{HTML}{8A5A00}
\definecolor{obeAmberLight}{HTML}{FDF1D6}
\definecolor{obeMaroon}{HTML}{7A2233}
\definecolor{obeMaroonLight}{HTML}{F7E1E5}
\definecolor{obeGray}{HTML}{4A4A4A}
\definecolor{obeGrayLight}{HTML}{F0F0F0}
\definecolor{codeBack}{HTML}{F7F7F4}
\definecolor{codeFrame}{HTML}{C9C9C0}

% ---------------------------------------------------------
% 3. Identitas mata kuliah (dipakai seluruh deck)
% ---------------------------------------------------------
\newcommand{\NamaMK}{Bahasa Pemrograman Tingkat Rendah}
\newcommand{\NamaMKShort}{Bhs.\ Pemrograman Tingkat Rendah}
\newcommand{\KodeMK}{MWCA010}
\newcommand{\Prodi}{Program Studi Teknik Informatika}
\newcommand{\Fakultas}{Fakultas Teknologi Informasi}
\newcommand{\PerguruanTinggi}{Universitas Bale Bandung}
\newcommand{\DosenPengampu}{Cecep Suwanda, S.Si., M.Kom.}
\newcommand{\SemesterTeks}{Semester Ganjil}

% \PertemuanTeks dan \JudulBab ditetapkan di masing-masing
% berkas minggu-NN.tex sebelum \begin{document}.
\providecommand{\PertemuanTeks}{Pertemuan}

% ---------------------------------------------------------
% 4. Tema beamer
% ---------------------------------------------------------
\usetheme{default}
\setbeamertemplate{navigation symbols}{}

\setbeamercolor{normal text}{fg=black!88}
\setbeamercolor{structure}{fg=obeBlue}
\setbeamercolor{alerted text}{fg=obeMaroon}

% --- judul slide: pita biru penuh lebar ---
\setbeamercolor{frametitle}{fg=white,bg=obeBlue}
\setbeamerfont{frametitle}{series=\bfseries,size=\large}
\setbeamertemplate{frametitle}{%
    \nointerlineskip%
    \begin{beamercolorbox}[%
        wd=\paperwidth,ht=2.9ex,dp=1.5ex,
        leftskip=0.4cm,rightskip=0.4cm]{frametitle}%
        \usebeamerfont{frametitle}\insertframetitle%
    \end{beamercolorbox}%
    \nointerlineskip%
    \hbox{\textcolor{obeTeal}{\rule{\paperwidth}{1.1pt}}}%
}

% --- kaki slide: identitas mata kuliah, pertemuan, nomor ---
\setbeamercolor{footline}{fg=white,bg=obeGray}
\setbeamerfont{footline}{size=\scriptsize,family=\sffamily}
\setbeamertemplate{footline}{%
    \leavevmode%
    \hbox{%
        \begin{beamercolorbox}[wd=0.42\paperwidth,ht=2.1ex,dp=1.1ex,
            leftskip=0.4cm]{footline}%
            \usebeamerfont{footline}%
            \KodeMK\ \textbullet\ \NamaMKShort%
        \end{beamercolorbox}%
        \begin{beamercolorbox}[wd=0.32\paperwidth,ht=2.1ex,dp=1.1ex,
            center]{footline}%
            \usebeamerfont{footline}\PertemuanTeks%
        \end{beamercolorbox}%
        \begin{beamercolorbox}[wd=0.26\paperwidth,ht=2.1ex,dp=1.1ex,
            rightskip=0.4cm]{footline}%
            \usebeamerfont{footline}%
            \insertframenumber\,/\,\inserttotalframenumber%
        \end{beamercolorbox}%
    }%
    \vskip0pt%
}

% --- daftar ---
\setbeamercolor{itemize item}{fg=obeTeal}
\setbeamercolor{itemize subitem}{fg=obeTeal}
\setbeamercolor{enumerate item}{fg=obeBlue}
\setbeamertemplate{itemize item}{\raise0.2ex\hbox{\tiny$\blacksquare$}}
\setbeamertemplate{itemize subitem}{\raise0.2ex\hbox{\scriptsize$\bullet$}}

% --- kotak (pengganti tcolorbox pada buku) ---
\setbeamercolor{block title}{bg=obeBlueLight,fg=obeBlue}
\setbeamercolor{block body}{bg=obeBlueLight!45,fg=black!88}
\setbeamercolor{block title alerted}{bg=obeMaroonLight,fg=obeMaroon}
\setbeamercolor{block body alerted}{bg=obeMaroonLight!45,fg=black!88}
\setbeamercolor{block title example}{bg=obeTealLight,fg=obeTeal}
\setbeamercolor{block body example}{bg=obeTealLight!45,fg=black!88}

% ---------------------------------------------------------
% 5. Makro notasi (sama persis dengan buku)
%
% Catatan: \note TIDAK didefinisikan di sini karena beamer
% sudah memakai \note untuk catatan pembicara. Untuk catatan
% di slide dipakai \catatan.
% ---------------------------------------------------------
\newcommand{\code}[1]{\texttt{#1}}
\newcommand{\asm}[1]{\texttt{#1}}
\newcommand{\reg}[1]{\texttt{#1}}
\newcommand{\flag}[1]{\texttt{#1}}
\newcommand{\dir}[1]{\texttt{#1}}
\newcommand{\file}[1]{\texttt{#1}}
\newcommand{\eng}[1]{\textit{#1}}

\newcommand{\hex}[1]{\texttt{#1h}}
\newcommand{\bin}[1]{\texttt{#1b}}
\newcommand{\dec}[1]{\texttt{#1}}

\newcommand{\important}[1]{\textbf{Penting:} #1}
\newcommand{\catatan}[1]{\textbf{Catatan:} #1}
\newcommand{\term}[1]{\textbf{#1}}

\newcommand{\CPL}[1]{\textbf{CPL#1}}
\newcommand{\CPMK}[1]{\textbf{CPMK#1}}
\newcommand{\SubCPMK}[1]{\textbf{Sub-CPMK #1}}

\newcommand{\checkbox}{\ding{113}\hspace{0.35em}}
\newcommand{\checked}{\ding{51}\hspace{0.35em}}
\newcommand{\crossed}{\ding{55}\hspace{0.35em}}

\newcommand{\flowarrow}{\ensuremath{\rightarrow}}

% ---------------------------------------------------------
% 6. Makro tata letak slide
% ---------------------------------------------------------

% Keterangan sumber materi (ketertelusuran ke buku)
\newcommand{\Sumber}[1]{%
    \vfill
    {\scriptsize\color{obeGray}Sumber: #1\par}%
}

% Frame judul deck: \JudulDeck{<judul bab>}{<pertemuan>}
\newcommand{\JudulDeck}[2]{%
\begin{frame}[plain]
    \vspace{0.25cm}
    {\color{obeBlue}\sffamily\bfseries\small \NamaMK\ (\KodeMK)\par}
    \vspace{0.12cm}
    {\color{obeGray}\sffamily\scriptsize
     \Prodi\ \textbullet\ \Fakultas\ \textbullet\ \PerguruanTinggi\par}
    \vfill
    {\color{obeBlue}\sffamily\LARGE\bfseries #1\par}
    \vspace{0.4cm}
    {\color{obeTeal}\sffamily\large #2\par}
    \vspace{0.45cm}
    {\color{obeTeal}\rule{\textwidth}{1pt}}
    \vfill
    {\color{obeGray}\sffamily\small Dosen Pengampu\par}
    \vspace{0.08cm}
    {\sffamily\bfseries \DosenPengampu\par}
    \vspace{0.2cm}
    {\color{obeGray}\sffamily\small \SemesterTeks\par}
    \vspace{0.25cm}
\end{frame}
}

% Frame pemisah bagian
\newcommand{\PemisahBagian}[1]{%
\begin{frame}[plain]
    \vfill
    \begin{center}
        {\color{obeTeal}\sffamily\small\bfseries BAGIAN\par}
        \vspace{0.35cm}
        {\color{obeBlue}\sffamily\Large\bfseries #1\par}
        \vspace{0.35cm}
        {\color{obeTeal}\rule{0.55\textwidth}{0.9pt}}
    \end{center}
    \vfill
\end{frame}
}

% ---------------------------------------------------------
% 7. Gaya penampil kode
% ---------------------------------------------------------
\input{listing-ppt}

%
% Palet warna disalin dari buku_ajar/config/styles.tex agar
% identitas visual slide sama dengan buku ajar.
% =========================================================

% ---------------------------------------------------------
% 1. Paket
% ---------------------------------------------------------
\usepackage[T1]{fontenc}
\usepackage[utf8]{inputenc}
\usepackage[indonesian]{babel}
\usepackage{lmodern}
\usepackage{tikz}
\usetikzlibrary{arrows.meta, positioning, calc, shapes.geometric}
\usepackage{booktabs}
\usepackage{array}
\usepackage{amsmath}
\usepackage{amsfonts}
\usepackage{pifont}          % \ding untuk tanda checklist
\usepackage{listings}

% ---------------------------------------------------------
% 2. Palet warna OBE (identik dengan buku)
% ---------------------------------------------------------
\definecolor{obeBlue}{HTML}{1F3864}
\definecolor{obeBlueLight}{HTML}{DCE4F5}
\definecolor{obeTeal}{HTML}{0F6E6E}
\definecolor{obeTealLight}{HTML}{D9F0EE}
\definecolor{obeGreen}{HTML}{2E6B34}
\definecolor{obeGreenLight}{HTML}{DEEFDF}
\definecolor{obeAmber}{HTML}{8A5A00}
\definecolor{obeAmberLight}{HTML}{FDF1D6}
\definecolor{obeMaroon}{HTML}{7A2233}
\definecolor{obeMaroonLight}{HTML}{F7E1E5}
\definecolor{obeGray}{HTML}{4A4A4A}
\definecolor{obeGrayLight}{HTML}{F0F0F0}
\definecolor{codeBack}{HTML}{F7F7F4}
\definecolor{codeFrame}{HTML}{C9C9C0}

% ---------------------------------------------------------
% 3. Identitas mata kuliah (dipakai seluruh deck)
% ---------------------------------------------------------
\newcommand{\NamaMK}{Bahasa Pemrograman Tingkat Rendah}
\newcommand{\NamaMKShort}{Bhs.\ Pemrograman Tingkat Rendah}
\newcommand{\KodeMK}{MWCA010}
\newcommand{\Prodi}{Program Studi Teknik Informatika}
\newcommand{\Fakultas}{Fakultas Teknologi Informasi}
\newcommand{\PerguruanTinggi}{Universitas Bale Bandung}
\newcommand{\DosenPengampu}{Cecep Suwanda, S.Si., M.Kom.}
\newcommand{\SemesterTeks}{Semester Ganjil}

% \PertemuanTeks dan \JudulBab ditetapkan di masing-masing
% berkas minggu-NN.tex sebelum \begin{document}.
\providecommand{\PertemuanTeks}{Pertemuan}

% ---------------------------------------------------------
% 4. Tema beamer
% ---------------------------------------------------------
\usetheme{default}
\setbeamertemplate{navigation symbols}{}

\setbeamercolor{normal text}{fg=black!88}
\setbeamercolor{structure}{fg=obeBlue}
\setbeamercolor{alerted text}{fg=obeMaroon}

% --- judul slide: pita biru penuh lebar ---
\setbeamercolor{frametitle}{fg=white,bg=obeBlue}
\setbeamerfont{frametitle}{series=\bfseries,size=\large}
\setbeamertemplate{frametitle}{%
    \nointerlineskip%
    \begin{beamercolorbox}[%
        wd=\paperwidth,ht=2.9ex,dp=1.5ex,
        leftskip=0.4cm,rightskip=0.4cm]{frametitle}%
        \usebeamerfont{frametitle}\insertframetitle%
    \end{beamercolorbox}%
    \nointerlineskip%
    \hbox{\textcolor{obeTeal}{\rule{\paperwidth}{1.1pt}}}%
}

% --- kaki slide: identitas mata kuliah, pertemuan, nomor ---
\setbeamercolor{footline}{fg=white,bg=obeGray}
\setbeamerfont{footline}{size=\scriptsize,family=\sffamily}
\setbeamertemplate{footline}{%
    \leavevmode%
    \hbox{%
        \begin{beamercolorbox}[wd=0.42\paperwidth,ht=2.1ex,dp=1.1ex,
            leftskip=0.4cm]{footline}%
            \usebeamerfont{footline}%
            \KodeMK\ \textbullet\ \NamaMKShort%
        \end{beamercolorbox}%
        \begin{beamercolorbox}[wd=0.32\paperwidth,ht=2.1ex,dp=1.1ex,
            center]{footline}%
            \usebeamerfont{footline}\PertemuanTeks%
        \end{beamercolorbox}%
        \begin{beamercolorbox}[wd=0.26\paperwidth,ht=2.1ex,dp=1.1ex,
            rightskip=0.4cm]{footline}%
            \usebeamerfont{footline}%
            \insertframenumber\,/\,\inserttotalframenumber%
        \end{beamercolorbox}%
    }%
    \vskip0pt%
}

% --- daftar ---
\setbeamercolor{itemize item}{fg=obeTeal}
\setbeamercolor{itemize subitem}{fg=obeTeal}
\setbeamercolor{enumerate item}{fg=obeBlue}
\setbeamertemplate{itemize item}{\raise0.2ex\hbox{\tiny$\blacksquare$}}
\setbeamertemplate{itemize subitem}{\raise0.2ex\hbox{\scriptsize$\bullet$}}

% --- kotak (pengganti tcolorbox pada buku) ---
\setbeamercolor{block title}{bg=obeBlueLight,fg=obeBlue}
\setbeamercolor{block body}{bg=obeBlueLight!45,fg=black!88}
\setbeamercolor{block title alerted}{bg=obeMaroonLight,fg=obeMaroon}
\setbeamercolor{block body alerted}{bg=obeMaroonLight!45,fg=black!88}
\setbeamercolor{block title example}{bg=obeTealLight,fg=obeTeal}
\setbeamercolor{block body example}{bg=obeTealLight!45,fg=black!88}

% ---------------------------------------------------------
% 5. Makro notasi (sama persis dengan buku)
%
% Catatan: \note TIDAK didefinisikan di sini karena beamer
% sudah memakai \note untuk catatan pembicara. Untuk catatan
% di slide dipakai \catatan.
% ---------------------------------------------------------
\newcommand{\code}[1]{\texttt{#1}}
\newcommand{\asm}[1]{\texttt{#1}}
\newcommand{\reg}[1]{\texttt{#1}}
\newcommand{\flag}[1]{\texttt{#1}}
\newcommand{\dir}[1]{\texttt{#1}}
\newcommand{\file}[1]{\texttt{#1}}
\newcommand{\eng}[1]{\textit{#1}}

\newcommand{\hex}[1]{\texttt{#1h}}
\newcommand{\bin}[1]{\texttt{#1b}}
\newcommand{\dec}[1]{\texttt{#1}}

\newcommand{\important}[1]{\textbf{Penting:} #1}
\newcommand{\catatan}[1]{\textbf{Catatan:} #1}
\newcommand{\term}[1]{\textbf{#1}}

\newcommand{\CPL}[1]{\textbf{CPL#1}}
\newcommand{\CPMK}[1]{\textbf{CPMK#1}}
\newcommand{\SubCPMK}[1]{\textbf{Sub-CPMK #1}}

\newcommand{\checkbox}{\ding{113}\hspace{0.35em}}
\newcommand{\checked}{\ding{51}\hspace{0.35em}}
\newcommand{\crossed}{\ding{55}\hspace{0.35em}}

\newcommand{\flowarrow}{\ensuremath{\rightarrow}}

% ---------------------------------------------------------
% 6. Makro tata letak slide
% ---------------------------------------------------------

% Keterangan sumber materi (ketertelusuran ke buku)
\newcommand{\Sumber}[1]{%
    \vfill
    {\scriptsize\color{obeGray}Sumber: #1\par}%
}

% Frame judul deck: \JudulDeck{<judul bab>}{<pertemuan>}
\newcommand{\JudulDeck}[2]{%
\begin{frame}[plain]
    \vspace{0.25cm}
    {\color{obeBlue}\sffamily\bfseries\small \NamaMK\ (\KodeMK)\par}
    \vspace{0.12cm}
    {\color{obeGray}\sffamily\scriptsize
     \Prodi\ \textbullet\ \Fakultas\ \textbullet\ \PerguruanTinggi\par}
    \vfill
    {\color{obeBlue}\sffamily\LARGE\bfseries #1\par}
    \vspace{0.4cm}
    {\color{obeTeal}\sffamily\large #2\par}
    \vspace{0.45cm}
    {\color{obeTeal}\rule{\textwidth}{1pt}}
    \vfill
    {\color{obeGray}\sffamily\small Dosen Pengampu\par}
    \vspace{0.08cm}
    {\sffamily\bfseries \DosenPengampu\par}
    \vspace{0.2cm}
    {\color{obeGray}\sffamily\small \SemesterTeks\par}
    \vspace{0.25cm}
\end{frame}
}

% Frame pemisah bagian
\newcommand{\PemisahBagian}[1]{%
\begin{frame}[plain]
    \vfill
    \begin{center}
        {\color{obeTeal}\sffamily\small\bfseries BAGIAN\par}
        \vspace{0.35cm}
        {\color{obeBlue}\sffamily\Large\bfseries #1\par}
        \vspace{0.35cm}
        {\color{obeTeal}\rule{0.55\textwidth}{0.9pt}}
    \end{center}
    \vfill
\end{frame}
}

% ---------------------------------------------------------
% 7. Gaya penampil kode
% ---------------------------------------------------------
% =========================================================
% ppt/src/listing-ppt.tex
% Gaya penampil kode untuk slide.
%
% Diturunkan dari buku_ajar/config/listings.tex, tetapi
% disesuaikan untuk slide: ukuran lebih kecil (karena frame
% 16:9 dan kode harus utuh terlihat) dan jarak lebih rapat.
% =========================================================

% ---------------------------------------------------------
% 1. Definisi bahasa MASM/TASM 8086 (salinan dari buku)
% ---------------------------------------------------------
\lstdefinelanguage{MASM}{
    morekeywords={
        MOV, XCHG, LEA, LDS, LES, PUSH, POP, PUSHA, POPA, PUSHF, POPF,
        LAHF, SAHF, XLAT, IN, OUT, CBW, CWD,
        ADD, ADC, SUB, SBB, INC, DEC, NEG, CMP, MUL, IMUL, DIV, IDIV,
        AAA, AAS, AAM, AAD, DAA, DAS,
        AND, OR, XOR, NOT, TEST, SHL, SHR, SAL, SAR, ROL, ROR, RCL, RCR,
        JMP, JE, JZ, JNE, JNZ, JG, JGE, JL, JLE, JA, JAE, JB, JBE,
        JC, JNC, JS, JNS, JO, JNO, JP, JPE, JNP, JPO, JCXZ,
        LOOP, LOOPE, LOOPZ, LOOPNE, LOOPNZ, CALL, RET, INT, IRET, INTO,
        HLT, NOP, CLI, STI, CLD, STD, CLC, STC, CMC, WAIT, ESC, LOCK,
        DB, DW, DD, DQ, DT, EQU, ORG, SEGMENT, ENDS, ASSUME, END,
        PROC, ENDP, MACRO, ENDM, LOCAL, PUBLIC, EXTRN, EXTERN,
        INCLUDE, INCLUDELIB, MODEL, STACK, CODE, DATA, CONST, ALIGN, EVEN,
        OFFSET, PTR, BYTE, WORD, DWORD, QWORD, TBYTE,
        NEAR, FAR, SHORT, LABEL, TYPE, SIZE, LENGTH, DUP, THIS,
        AX, BX, CX, DX, AL, AH, BL, BH, CL, CH, DL, DH,
        SP, BP, SI, DI, IP, CS, DS, SS, ES, FLAGS,
        CF, PF, AF, ZF, TF, IF, DF, OF
    },
    sensitive=false,
    morecomment=[l]{;},
    morestring=[b]',
    morestring=[b]"
}

% ---------------------------------------------------------
% 2. Gaya kode assembly untuk slide
% ---------------------------------------------------------
\lstdefinestyle{asmslide}{
    language=MASM,
    basicstyle=\ttfamily\scriptsize,
    keywordstyle=\color{obeBlue}\bfseries,
    commentstyle=\color{obeGreen}\itshape,
    stringstyle=\color{obeMaroon},
    numberstyle=\tiny\color{obeGray},
    numbers=left,
    numbersep=5pt,
    stepnumber=1,
    backgroundcolor=\color{codeBack},
    frame=single,
    rulecolor=\color{codeFrame},
    framesep=4pt,
    breaklines=true,
    breakatwhitespace=false,
    % CATATAN PENTING: opsi `postbreak` TIDAK boleh dipakai di sini.
    % Di beamer, `breaklines=true` bersama `postbreak=<kode>` membuat
    % listings gagal dirakit dengan galat
    %   "Missing number, treated as zero <to be read again> \kern".
    % Karena \lstset{style=asmslide} bersifat global, opsi postbreak
    % juga bocor ke gaya lain (mis. outslide/keluaran) dan merusaknya.
    % Di buku ajar (kelas `book`) justru sebaliknya: postbreak aman dan
    % dipakai. Jadi larangan ini khusus untuk slide.
    showstringspaces=false,
    tabsize=4,
    columns=fullflexible,
    keepspaces=true,
    upquote=true,
    xleftmargin=12pt,
    aboveskip=0.4em,
    belowskip=0.4em
}

% ---------------------------------------------------------
% 3. Gaya keluaran program / sesi debugger untuk slide
% ---------------------------------------------------------
\lstdefinestyle{outslide}{
    basicstyle=\ttfamily\scriptsize,
    keywordstyle=\color{black},
    commentstyle=\color{obeGray},
    stringstyle=\color{black},
    numbers=none,
    backgroundcolor=\color{obeGrayLight},
    frame=single,
    rulecolor=\color{codeFrame},
    framesep=4pt,
    breaklines=true,
    showstringspaces=false,
    columns=fullflexible,
    keepspaces=true,
    upquote=true,
    xleftmargin=4pt,
    aboveskip=0.4em,
    belowskip=0.4em
}

\lstset{style=asmslide}

% ---------------------------------------------------------
% 4. Environment ringkas untuk dipakai di dalam frame
%
% Frame yang memuat kedua environment ini WAJIB memakai
% opsi [fragile], karena isinya verbatim.
% ---------------------------------------------------------
\lstnewenvironment{kodeasm}[1][]{%
    \lstset{style=asmslide, #1}%
}{}

\lstnewenvironment{keluaran}[1][]{%
    \lstset{style=outslide, #1}%
}{}


%
% Palet warna disalin dari buku_ajar/config/styles.tex agar
% identitas visual slide sama dengan buku ajar.
% =========================================================

% ---------------------------------------------------------
% 1. Paket
% ---------------------------------------------------------
\usepackage[T1]{fontenc}
\usepackage[utf8]{inputenc}
\usepackage[indonesian]{babel}
\usepackage{lmodern}
\usepackage{tikz}
\usetikzlibrary{arrows.meta, positioning, calc, shapes.geometric}
\usepackage{booktabs}
\usepackage{array}
\usepackage{amsmath}
\usepackage{amsfonts}
\usepackage{pifont}          % \ding untuk tanda checklist
\usepackage{listings}

% ---------------------------------------------------------
% 2. Palet warna OBE (identik dengan buku)
% ---------------------------------------------------------
\definecolor{obeBlue}{HTML}{1F3864}
\definecolor{obeBlueLight}{HTML}{DCE4F5}
\definecolor{obeTeal}{HTML}{0F6E6E}
\definecolor{obeTealLight}{HTML}{D9F0EE}
\definecolor{obeGreen}{HTML}{2E6B34}
\definecolor{obeGreenLight}{HTML}{DEEFDF}
\definecolor{obeAmber}{HTML}{8A5A00}
\definecolor{obeAmberLight}{HTML}{FDF1D6}
\definecolor{obeMaroon}{HTML}{7A2233}
\definecolor{obeMaroonLight}{HTML}{F7E1E5}
\definecolor{obeGray}{HTML}{4A4A4A}
\definecolor{obeGrayLight}{HTML}{F0F0F0}
\definecolor{codeBack}{HTML}{F7F7F4}
\definecolor{codeFrame}{HTML}{C9C9C0}

% ---------------------------------------------------------
% 3. Identitas mata kuliah (dipakai seluruh deck)
% ---------------------------------------------------------
\newcommand{\NamaMK}{Bahasa Pemrograman Tingkat Rendah}
\newcommand{\NamaMKShort}{Bhs.\ Pemrograman Tingkat Rendah}
\newcommand{\KodeMK}{MWCA010}
\newcommand{\Prodi}{Program Studi Teknik Informatika}
\newcommand{\Fakultas}{Fakultas Teknologi Informasi}
\newcommand{\PerguruanTinggi}{Universitas Bale Bandung}
\newcommand{\DosenPengampu}{Cecep Suwanda, S.Si., M.Kom.}
\newcommand{\SemesterTeks}{Semester Ganjil}

% \PertemuanTeks dan \JudulBab ditetapkan di masing-masing
% berkas minggu-NN.tex sebelum \begin{document}.
\providecommand{\PertemuanTeks}{Pertemuan}

% ---------------------------------------------------------
% 4. Tema beamer
% ---------------------------------------------------------
\usetheme{default}
\setbeamertemplate{navigation symbols}{}

\setbeamercolor{normal text}{fg=black!88}
\setbeamercolor{structure}{fg=obeBlue}
\setbeamercolor{alerted text}{fg=obeMaroon}

% --- judul slide: pita biru penuh lebar ---
\setbeamercolor{frametitle}{fg=white,bg=obeBlue}
\setbeamerfont{frametitle}{series=\bfseries,size=\large}
\setbeamertemplate{frametitle}{%
    \nointerlineskip%
    \begin{beamercolorbox}[%
        wd=\paperwidth,ht=2.9ex,dp=1.5ex,
        leftskip=0.4cm,rightskip=0.4cm]{frametitle}%
        \usebeamerfont{frametitle}\insertframetitle%
    \end{beamercolorbox}%
    \nointerlineskip%
    \hbox{\textcolor{obeTeal}{\rule{\paperwidth}{1.1pt}}}%
}

% --- kaki slide: identitas mata kuliah, pertemuan, nomor ---
\setbeamercolor{footline}{fg=white,bg=obeGray}
\setbeamerfont{footline}{size=\scriptsize,family=\sffamily}
\setbeamertemplate{footline}{%
    \leavevmode%
    \hbox{%
        \begin{beamercolorbox}[wd=0.42\paperwidth,ht=2.1ex,dp=1.1ex,
            leftskip=0.4cm]{footline}%
            \usebeamerfont{footline}%
            \KodeMK\ \textbullet\ \NamaMKShort%
        \end{beamercolorbox}%
        \begin{beamercolorbox}[wd=0.32\paperwidth,ht=2.1ex,dp=1.1ex,
            center]{footline}%
            \usebeamerfont{footline}\PertemuanTeks%
        \end{beamercolorbox}%
        \begin{beamercolorbox}[wd=0.26\paperwidth,ht=2.1ex,dp=1.1ex,
            rightskip=0.4cm]{footline}%
            \usebeamerfont{footline}%
            \insertframenumber\,/\,\inserttotalframenumber%
        \end{beamercolorbox}%
    }%
    \vskip0pt%
}

% --- daftar ---
\setbeamercolor{itemize item}{fg=obeTeal}
\setbeamercolor{itemize subitem}{fg=obeTeal}
\setbeamercolor{enumerate item}{fg=obeBlue}
\setbeamertemplate{itemize item}{\raise0.2ex\hbox{\tiny$\blacksquare$}}
\setbeamertemplate{itemize subitem}{\raise0.2ex\hbox{\scriptsize$\bullet$}}

% --- kotak (pengganti tcolorbox pada buku) ---
\setbeamercolor{block title}{bg=obeBlueLight,fg=obeBlue}
\setbeamercolor{block body}{bg=obeBlueLight!45,fg=black!88}
\setbeamercolor{block title alerted}{bg=obeMaroonLight,fg=obeMaroon}
\setbeamercolor{block body alerted}{bg=obeMaroonLight!45,fg=black!88}
\setbeamercolor{block title example}{bg=obeTealLight,fg=obeTeal}
\setbeamercolor{block body example}{bg=obeTealLight!45,fg=black!88}

% ---------------------------------------------------------
% 5. Makro notasi (sama persis dengan buku)
%
% Catatan: \note TIDAK didefinisikan di sini karena beamer
% sudah memakai \note untuk catatan pembicara. Untuk catatan
% di slide dipakai \catatan.
% ---------------------------------------------------------
\newcommand{\code}[1]{\texttt{#1}}
\newcommand{\asm}[1]{\texttt{#1}}
\newcommand{\reg}[1]{\texttt{#1}}
\newcommand{\flag}[1]{\texttt{#1}}
\newcommand{\dir}[1]{\texttt{#1}}
\newcommand{\file}[1]{\texttt{#1}}
\newcommand{\eng}[1]{\textit{#1}}

\newcommand{\hex}[1]{\texttt{#1h}}
\newcommand{\bin}[1]{\texttt{#1b}}
\newcommand{\dec}[1]{\texttt{#1}}

\newcommand{\important}[1]{\textbf{Penting:} #1}
\newcommand{\catatan}[1]{\textbf{Catatan:} #1}
\newcommand{\term}[1]{\textbf{#1}}

\newcommand{\CPL}[1]{\textbf{CPL#1}}
\newcommand{\CPMK}[1]{\textbf{CPMK#1}}
\newcommand{\SubCPMK}[1]{\textbf{Sub-CPMK #1}}

\newcommand{\checkbox}{\ding{113}\hspace{0.35em}}
\newcommand{\checked}{\ding{51}\hspace{0.35em}}
\newcommand{\crossed}{\ding{55}\hspace{0.35em}}

\newcommand{\flowarrow}{\ensuremath{\rightarrow}}

% ---------------------------------------------------------
% 6. Makro tata letak slide
% ---------------------------------------------------------

% Keterangan sumber materi (ketertelusuran ke buku)
\newcommand{\Sumber}[1]{%
    \vfill
    {\scriptsize\color{obeGray}Sumber: #1\par}%
}

% Frame judul deck: \JudulDeck{<judul bab>}{<pertemuan>}
\newcommand{\JudulDeck}[2]{%
\begin{frame}[plain]
    \vspace{0.25cm}
    {\color{obeBlue}\sffamily\bfseries\small \NamaMK\ (\KodeMK)\par}
    \vspace{0.12cm}
    {\color{obeGray}\sffamily\scriptsize
     \Prodi\ \textbullet\ \Fakultas\ \textbullet\ \PerguruanTinggi\par}
    \vfill
    {\color{obeBlue}\sffamily\LARGE\bfseries #1\par}
    \vspace{0.4cm}
    {\color{obeTeal}\sffamily\large #2\par}
    \vspace{0.45cm}
    {\color{obeTeal}\rule{\textwidth}{1pt}}
    \vfill
    {\color{obeGray}\sffamily\small Dosen Pengampu\par}
    \vspace{0.08cm}
    {\sffamily\bfseries \DosenPengampu\par}
    \vspace{0.2cm}
    {\color{obeGray}\sffamily\small \SemesterTeks\par}
    \vspace{0.25cm}
\end{frame}
}

% Frame pemisah bagian
\newcommand{\PemisahBagian}[1]{%
\begin{frame}[plain]
    \vfill
    \begin{center}
        {\color{obeTeal}\sffamily\small\bfseries BAGIAN\par}
        \vspace{0.35cm}
        {\color{obeBlue}\sffamily\Large\bfseries #1\par}
        \vspace{0.35cm}
        {\color{obeTeal}\rule{0.55\textwidth}{0.9pt}}
    \end{center}
    \vfill
\end{frame}
}

% ---------------------------------------------------------
% 7. Gaya penampil kode
% ---------------------------------------------------------
% =========================================================
% ppt/src/listing-ppt.tex
% Gaya penampil kode untuk slide.
%
% Diturunkan dari buku_ajar/config/listings.tex, tetapi
% disesuaikan untuk slide: ukuran lebih kecil (karena frame
% 16:9 dan kode harus utuh terlihat) dan jarak lebih rapat.
% =========================================================

% ---------------------------------------------------------
% 1. Definisi bahasa MASM/TASM 8086 (salinan dari buku)
% ---------------------------------------------------------
\lstdefinelanguage{MASM}{
    morekeywords={
        MOV, XCHG, LEA, LDS, LES, PUSH, POP, PUSHA, POPA, PUSHF, POPF,
        LAHF, SAHF, XLAT, IN, OUT, CBW, CWD,
        ADD, ADC, SUB, SBB, INC, DEC, NEG, CMP, MUL, IMUL, DIV, IDIV,
        AAA, AAS, AAM, AAD, DAA, DAS,
        AND, OR, XOR, NOT, TEST, SHL, SHR, SAL, SAR, ROL, ROR, RCL, RCR,
        JMP, JE, JZ, JNE, JNZ, JG, JGE, JL, JLE, JA, JAE, JB, JBE,
        JC, JNC, JS, JNS, JO, JNO, JP, JPE, JNP, JPO, JCXZ,
        LOOP, LOOPE, LOOPZ, LOOPNE, LOOPNZ, CALL, RET, INT, IRET, INTO,
        HLT, NOP, CLI, STI, CLD, STD, CLC, STC, CMC, WAIT, ESC, LOCK,
        DB, DW, DD, DQ, DT, EQU, ORG, SEGMENT, ENDS, ASSUME, END,
        PROC, ENDP, MACRO, ENDM, LOCAL, PUBLIC, EXTRN, EXTERN,
        INCLUDE, INCLUDELIB, MODEL, STACK, CODE, DATA, CONST, ALIGN, EVEN,
        OFFSET, PTR, BYTE, WORD, DWORD, QWORD, TBYTE,
        NEAR, FAR, SHORT, LABEL, TYPE, SIZE, LENGTH, DUP, THIS,
        AX, BX, CX, DX, AL, AH, BL, BH, CL, CH, DL, DH,
        SP, BP, SI, DI, IP, CS, DS, SS, ES, FLAGS,
        CF, PF, AF, ZF, TF, IF, DF, OF
    },
    sensitive=false,
    morecomment=[l]{;},
    morestring=[b]',
    morestring=[b]"
}

% ---------------------------------------------------------
% 2. Gaya kode assembly untuk slide
% ---------------------------------------------------------
\lstdefinestyle{asmslide}{
    language=MASM,
    basicstyle=\ttfamily\scriptsize,
    keywordstyle=\color{obeBlue}\bfseries,
    commentstyle=\color{obeGreen}\itshape,
    stringstyle=\color{obeMaroon},
    numberstyle=\tiny\color{obeGray},
    numbers=left,
    numbersep=5pt,
    stepnumber=1,
    backgroundcolor=\color{codeBack},
    frame=single,
    rulecolor=\color{codeFrame},
    framesep=4pt,
    breaklines=true,
    breakatwhitespace=false,
    % CATATAN PENTING: opsi `postbreak` TIDAK boleh dipakai di sini.
    % Di beamer, `breaklines=true` bersama `postbreak=<kode>` membuat
    % listings gagal dirakit dengan galat
    %   "Missing number, treated as zero <to be read again> \kern".
    % Karena \lstset{style=asmslide} bersifat global, opsi postbreak
    % juga bocor ke gaya lain (mis. outslide/keluaran) dan merusaknya.
    % Di buku ajar (kelas `book`) justru sebaliknya: postbreak aman dan
    % dipakai. Jadi larangan ini khusus untuk slide.
    showstringspaces=false,
    tabsize=4,
    columns=fullflexible,
    keepspaces=true,
    upquote=true,
    xleftmargin=12pt,
    aboveskip=0.4em,
    belowskip=0.4em
}

% ---------------------------------------------------------
% 3. Gaya keluaran program / sesi debugger untuk slide
% ---------------------------------------------------------
\lstdefinestyle{outslide}{
    basicstyle=\ttfamily\scriptsize,
    keywordstyle=\color{black},
    commentstyle=\color{obeGray},
    stringstyle=\color{black},
    numbers=none,
    backgroundcolor=\color{obeGrayLight},
    frame=single,
    rulecolor=\color{codeFrame},
    framesep=4pt,
    breaklines=true,
    showstringspaces=false,
    columns=fullflexible,
    keepspaces=true,
    upquote=true,
    xleftmargin=4pt,
    aboveskip=0.4em,
    belowskip=0.4em
}

\lstset{style=asmslide}

% ---------------------------------------------------------
% 4. Environment ringkas untuk dipakai di dalam frame
%
% Frame yang memuat kedua environment ini WAJIB memakai
% opsi [fragile], karena isinya verbatim.
% ---------------------------------------------------------
\lstnewenvironment{kodeasm}[1][]{%
    \lstset{style=asmslide, #1}%
}{}

\lstnewenvironment{keluaran}[1][]{%
    \lstset{style=outslide, #1}%
}{}



\renewcommand{\PertemuanTeks}{Pertemuan 3}

\begin{document}

% =========================================================
% 1. JUDUL
% =========================================================
\JudulDeck{Register 8086 dan Register Flag}{Pertemuan 3}

% =========================================================
% 2. CAPAIAN DAN TUJUAN
% =========================================================
\begin{frame}{Capaian dan Tujuan Pembelajaran}
\begin{block}{Sub-CPMK 1.2}
\small Mengidentifikasi fungsi register umum (\reg{AX}, \reg{BX},
\reg{CX}, \reg{DX}), register segmen (\reg{CS}, \reg{DS}, \reg{SS},
\reg{ES}), register \eng{pointer}/\eng{index} (\reg{SP}, \reg{BP},
\reg{SI}, \reg{DI}), dan register \eng{flag} beserta pengaruhnya
terhadap eksekusi program.
\end{block}

\vspace{0.15em}
{\small\textbf{CPMK yang didukung:} \CPMK{01}\quad\textbullet\quad
\textbf{CPL yang didukung:} \CPL{02}\par}

\vspace{0.25em}
\begin{exampleblock}{Setelah pertemuan ini, mahasiswa mampu:}
\footnotesize
\begin{enumerate}
    \item menyebutkan register umum 8086 beserta akses byte rendah
          dan byte tingginya;
    \item menjelaskan peran khas setiap register umum dan alasan di
          balik peran tersebut;
    \item menjelaskan fungsi register segmen serta peran register
          \eng{pointer}, \eng{index}, dan \reg{IP};
    \item mengidentifikasi sembilan \eng{flag} pada register
          \reg{FLAGS} beserta posisi bitnya;
    \item memprediksi perubahan \eng{flag} akibat operasi aritmatika
          dan menjelaskan akibatnya bagi alur program.
\end{enumerate}
\end{exampleblock}
\end{frame}

% =========================================================
% 3. PETA MATERI
% =========================================================
\begin{frame}{Peta Materi}
\small
\begin{enumerate}
    \item \textbf{Register Umum dan Register Segmen}\\
          {\footnotesize \reg{AX}--\reg{DX} dan akses byte
          rendah/tinggi; peran khas tiap register; \reg{CS},
          \reg{DS}, \reg{SS}, \reg{ES}}

    \vspace{0.25em}
    \item \textbf{Register \eng{Pointer}, \eng{Index}, dan
          \eng{Instruction Pointer}}\\
          {\footnotesize \reg{SP} dan \reg{BP} pada \eng{stack};
          \reg{SI} dan \reg{DI} pada operasi string; \reg{IP};
          register yang rawan diubah}

    \vspace{0.25em}
    \item \textbf{Register Flag}\\
          {\footnotesize peta bit \reg{FLAGS}; sembilan \eng{flag}
          status dan kendali; contoh perubahan \eng{flag}; lompatan
          bersyarat}
\end{enumerate}

\vspace{0.4em}
\begin{block}{Praktik pada pertemuan ini}
\small Menjalankan tiga contoh program, memprediksi \eng{flag}, lalu
mengamati register melalui \eng{Single Step}.
\end{block}
\end{frame}

% =========================================================
% BAGIAN 1
% =========================================================
\PemisahBagian{Register Umum dan Register Segmen}

% ---------------------------------------------------------
\begin{frame}{Empat Belas Register, Empat Kelompok}
\begin{center}
\begin{tikzpicture}[
    reg/.style={
        draw=obeBlue, line width=0.7pt, fill=obeBlueLight,
        rounded corners=2pt, minimum width=1.05cm, minimum height=0.55cm,
        align=center, font=\small\ttfamily, inner sep=2pt
    },
    regk/.style={reg, fill=obeTealLight, draw=obeTeal},
    regs/.style={reg, fill=obeAmberLight, draw=obeAmber},
    regf/.style={reg, fill=obeGrayLight, draw=obeGray},
    judul/.style={
        font=\scriptsize\bfseries\sffamily, text=obeBlue,
        anchor=east, align=right
    },
    ket/.style={font=\tiny\itshape, text=obeGray}
]
    \node[judul] at (-0.30, 4.55) {Register umum};
    \node[reg] (ax) at (0.35, 4.55) {AX\\{\normalfont\tiny akumulator}};
    \node[reg, right=0.18cm of ax] (bx) {BX\\{\normalfont\tiny basis}};
    \node[reg, right=0.18cm of bx] (cx) {CX\\{\normalfont\tiny pencacah}};
    \node[reg, right=0.18cm of cx] (dx) {DX\\{\normalfont\tiny data}};

    \node[judul] at (-0.30, 3.25) {Register segmen};
    \node[regk] (cs) at (0.35, 3.25) {CS\\{\normalfont\tiny kode}};
    \node[regk, right=0.18cm of cs] (ds) {DS\\{\normalfont\tiny data}};
    \node[regk, right=0.18cm of ds] (ss) {SS\\{\normalfont\tiny stack}};
    \node[regk, right=0.18cm of ss] (es) {ES\\{\normalfont\tiny tambahan}};

    \node[judul] at (-0.30, 1.95) {Pointer, index,\\dan IP};
    \node[regs] (sp) at (0.35, 1.95)
        {SP\\{\normalfont\tiny puncak stack}};
    \node[regs, right=0.18cm of sp] (bp)
        {BP\\{\normalfont\tiny acuan stack}};
    \node[regs, right=0.18cm of bp] (si)
        {SI\\{\normalfont\tiny sumber}};
    \node[regs, right=0.18cm of si] (di)
        {DI\\{\normalfont\tiny tujuan}};
    \node[regs, right=0.18cm of di] (ip)
        {IP\\{\normalfont\tiny instruksi}};

    \node[judul] at (-0.30, 0.60) {Register flag};
    \node[regf, minimum width=5.9cm] (fl) at (3.30, 0.60)
        {\textbf{FLAGS} (16 bit)\\{\normalfont\tiny
         CF PF AF ZF SF TF IF DF OF}};
\end{tikzpicture}
\end{center}

\vspace{0.3em}
{\footnotesize\color{obeGray}Keempat belas register berukuran 16 bit
dan dapat menampung nilai \hex{0000} sampai \hex{FFFF}. Register
adalah memori tercepat di dalam prosesor: satu daur mesin, tanpa
melewati bus.\par}
\Sumber{Bab V, Ringkasan kelompok register 8086}
\end{frame}

% ---------------------------------------------------------
\begin{frame}{Register Umum dan Akses Byte}
\small
\begin{center}
\begin{tabular}{@{}lllp{7.3cm}@{}}
    \toprule
    \textbf{Register} & \textbf{Byte tinggi} & \textbf{Byte rendah} &
    \textbf{Peran khas} \\
    \midrule
    \reg{AX} & \reg{AH} & \reg{AL} &
    Akumulator: hasil hitungan, operand wajib \asm{MUL}/\asm{DIV},
    layanan DOS \\
    \addlinespace
    \reg{BX} & \reg{BH} & \reg{BL} &
    Basis: penunjuk alamat memori pada mode terindeks \\
    \addlinespace
    \reg{CX} & \reg{CH} & \reg{CL} &
    Pencacah: pengulangan \asm{LOOP}, operasi string, jumlah
    pergeseran \asm{SHL}/\asm{SHR} \\
    \addlinespace
    \reg{DX} & \reg{DH} & \reg{DL} &
    Data: pasangan \reg{AX} pada operasi 32 bit, alamat port I/O dan
    layanan DOS \\
    \bottomrule
\end{tabular}
\end{center}

\vspace{0.3em}
{\footnotesize\color{obeGray}Setiap register 16 bit terbagi menjadi
\emph{byte tinggi} dan \emph{byte rendah} yang dapat diakses
sendiri-sendiri.\par}
\Sumber{Bab V, Register Umum dan Akses Byte}
\end{frame}

% ---------------------------------------------------------
\begin{frame}[fragile]{Akses Byte: \code{MOV AL, 5} dan
\code{MOV AX, 5}}
\small Dua instruksi ini terlihat serupa, tetapi akibatnya berbeda
pada \reg{AH}.
\vspace{0.3em}
\begin{columns}[t]
\begin{column}{0.54\textwidth}
\begin{kodeasm}[numbers=none]
MOV AL, 5   ; hanya byte rendah yang berubah
MOV AX, 5   ; kedua bagian ditulis sekaligus
\end{kodeasm}
\end{column}
\begin{column}{0.42\textwidth}
{\color{obeTeal}\bfseries Isi \reg{AX}}\\[0.3em]
\begin{keluaran}
sebelum       : FF00h
MOV AL, 5     : FF05h
MOV AX, 5     : 0005h
\end{keluaran}
\end{column}
\end{columns}

\vspace{0.4em}
\begin{alertblock}{Yang sering terlewat}
\small \code{MOV AL, 5} \textbf{tidak} mengosongkan \reg{AH}. Bila
\reg{AH} semula \hex{FF}, maka \reg{AX} menjadi \hex{FF05}, bukan
\hex{0005}.
\end{alertblock}
\end{frame}

% ---------------------------------------------------------
\begin{frame}{Mengapa Peran Setiap Register Berbeda}
\small Keempat register umum sama-sama dapat menampung nilai apa pun,
tetapi peran khasnya berbeda karena cara prosesor 8086 menyandikan
instruksi.

\vspace{0.4em}
\begin{center}
\begin{tabular}{@{}llp{7.6cm}@{}}
    \toprule
    \textbf{Register} & \textbf{Peran} & \textbf{Mengapa demikian} \\
    \midrule
    \reg{AX} & Akumulator &
    \asm{MUL}, \asm{DIV}, \asm{IN}, \asm{OUT}, dan sebagian besar
    layanan DOS mengandaikan operand di \reg{AX}/\reg{AL}; bentuk
    mesinnya juga lebih pendek \\
    \addlinespace
    \reg{BX} & Basis &
    Hanya \reg{BX} yang boleh menjadi register basis di dalam tanda
    kurung: \code{MOV AL, [BX]} \\
    \addlinespace
    \reg{CX} & Pencacah &
    \asm{LOOP}, \asm{REP}, dan operasi string memakainya sebagai
    penghitung otomatis; \asm{SHL}/\asm{SHR} memakai \reg{CL} \\
    \addlinespace
    \reg{DX} & Data &
    Pasangan \reg{AX} pada operasi 32 bit --- \reg{DX} menyimpan
    bagian tinggi --- serta alamat string dan nomor port \\
    \bottomrule
\end{tabular}
\end{center}

\vspace{0.3em}
{\footnotesize\color{obeGray}Register umum boleh dipakai di luar
peran khasnya; yang tidak boleh adalah melanggar aturan khusus,
misalnya \code{MOV AX, [CX]}.\par}
\end{frame}

% ---------------------------------------------------------
\begin{frame}{Register Segmen}
\small
\begin{center}
\begin{tabular}{@{}llp{7.6cm}@{}}
    \toprule
    \textbf{Register} & \textbf{Nama lengkap} & \textbf{Peran} \\
    \midrule
    \reg{CS} & \eng{Code Segment} &
    Menunjuk segmen tempat instruksi berada; selalu berpasangan
    dengan \reg{IP} \\
    \addlinespace
    \reg{DS} & \eng{Data Segment} &
    Menunjuk segmen tempat variabel berada; dipakai pada hampir
    semua akses data \\
    \addlinespace
    \reg{SS} & \eng{Stack Segment} &
    Menunjuk segmen tempat \eng{stack} berada; berpasangan dengan
    \reg{SP} dan \reg{BP} \\
    \addlinespace
    \reg{ES} & \eng{Extra Segment} &
    Segmen tambahan; menjadi segmen tujuan operasi string bersama
    \reg{DI} \\
    \bottomrule
\end{tabular}
\end{center}

\vspace{0.3em}
\begin{block}{Yang disimpan register segmen}
\small Bukan data yang dihitung, melainkan \emph{nomor segmen}
tempat data atau instruksi berada.
\end{block}
\Sumber{Bab V, Register Segmen}
\end{frame}

% ---------------------------------------------------------
\begin{frame}{Pasangan Segmen--Offset dan Inisialisasi \reg{DS}}
\small Setiap akses memori bermain dengan pasangan register
segmen--offset:
\vspace{0.3em}
\begin{columns}[t]
\begin{column}{0.48\textwidth}
\begin{itemize}
    \item \reg{CS}:\reg{IP} --- alamat instruksi berikutnya.
    \item \reg{DS}:\reg{SI} atau \reg{DS}:\reg{BX} --- alamat data
          umum.
    \item \reg{SS}:\reg{SP} atau \reg{SS}:\reg{BP} --- alamat di
          dalam \eng{stack}.
    \item \reg{ES}:\reg{DI} --- alamat tujuan operasi string.
\end{itemize}
\end{column}
\begin{column}{0.48\textwidth}
\begin{alertblock}{Dua kesalahan umum}
\footnotesize
\begin{enumerate}
    \item \code{MOV DS, 1000h} \textbf{tidak sah}. Nilai hanya boleh
          dari register umum: \code{MOV AX, 1000h} lalu
          \code{MOV DS, AX}.
    \item Mengubah \reg{DS} tanpa menyimpan nilai lamanya, sehingga
          variabel pada segmen data lama tidak terjangkau.
\end{enumerate}
\end{alertblock}
\end{column}
\end{columns}

\vspace{0.3em}
{\footnotesize\color{obeGray}Karena itu program \file{.EXE} wajib
menginisialisasi \reg{DS} di awal program.\par}
\end{frame}

% ---------------------------------------------------------
\begin{frame}{Rangkuman Bagian 1}
\small
\begin{itemize}
    \item Prosesor 8086 memiliki empat belas register 16 bit:
          empat register umum, empat register segmen, empat register
          \eng{pointer}/\eng{index}, ditambah \reg{IP} dan
          \reg{FLAGS}.

    \vspace{0.25em}
    \item \reg{AX}, \reg{BX}, \reg{CX}, dan \reg{DX} dapat diakses
          sebagai dua byte terpisah. Menulis ke byte rendah
          \textbf{tidak} mengubah byte tinggi.

    \vspace{0.25em}
    \item Peran khas muncul dari cara instruksi disandikan: \reg{AX}
          akumulator, \reg{BX} basis pengalamatan, \reg{CX} pencacah,
          \reg{DX} pasangan data dan penunjuk port.

    \vspace{0.25em}
    \item Register segmen menyimpan nomor segmen. Pasangan bakunya:
          \reg{CS}:\reg{IP}, \reg{DS}:\reg{SI}, \reg{SS}:\reg{SP},
          \reg{ES}:\reg{DI}.
\end{itemize}
\end{frame}

% =========================================================
% BAGIAN 2
% =========================================================
\PemisahBagian{Register Pointer, Index, dan Instruction Pointer}

% ---------------------------------------------------------
\begin{frame}{Stack Pointer dan Base Pointer}
\small
\begin{description}
    \item[\reg{SP} (\eng{Stack Pointer})] Selalu menunjuk puncak
          \eng{stack} saat ini. Setiap \asm{PUSH} mengurangi
          \reg{SP} sebesar 2 dan setiap \asm{POP} menambahnya
          sebesar 2, karena satu kata berukuran 2 byte. Nilainya
          berpasangan dengan \reg{SS}.

    \vspace{0.4em}
    \item[\reg{BP} (\eng{Base Pointer})] Menjadi titik acuan tetap
          di dalam \eng{stack}. Berbeda dengan \reg{SP} yang terus
          bergerak, \reg{BP} disetel sekali di awal prosedur lalu
          dipakai untuk menjangkau parameter dan variabel lokal.
\end{description}

\vspace{0.5em}
\begin{alertblock}{Akibat mengubah \reg{SP} sembarangan}
\small Puncak \eng{stack} ikut berpindah, sehingga instruksi
\asm{RET} dapat melompat ke alamat yang salah.
\end{alertblock}
\Sumber{Bab V, Stack Pointer dan Base Pointer}
\end{frame}

% ---------------------------------------------------------
\begin{frame}{Source Index dan Destination Index}
\small Dua register ini dipakai untuk mengakses deretan data, misalnya
\eng{string} atau \eng{array}.

\vspace{0.4em}
\begin{columns}[t]
\begin{column}{0.48\textwidth}
\begin{exampleblock}{\reg{SI} --- \eng{Source Index}}
\small Menunjuk alamat \emph{sumber} pada operasi string.
Pasangannya \reg{DS}, sehingga alamatnya ditulis
\reg{DS}:\reg{SI}.
\end{exampleblock}
\end{column}
\begin{column}{0.48\textwidth}
\begin{exampleblock}{\reg{DI} --- \eng{Destination Index}}
\small Menunjuk alamat \emph{tujuan} pada operasi string.
Pasangannya \reg{ES}, sehingga alamatnya ditulis
\reg{ES}:\reg{DI}.
\end{exampleblock}
\end{column}
\end{columns}

\vspace{0.5em}
\begin{block}{Memindahkan data antar segmen}
\small \asm{MOVSB}, \asm{MOVSW}, dan \asm{REP MOVSB} memindahkan data
dari \reg{DS}:\reg{SI} ke \reg{ES}:\reg{DI} sambil menyesuaikan
kedua register secara otomatis --- karena itu keduanya harus diisi
lebih dahulu.
\end{block}
\end{frame}

% ---------------------------------------------------------
\begin{frame}{Instruction Pointer dan Register yang Rawan}
\small
\begin{exampleblock}{\reg{IP} --- \eng{Instruction Pointer}}
\small
\begin{itemize}
    \item Menyimpan \eng{offset} instruksi berikutnya yang akan
          diambil; alamat lengkapnya ditentukan pasangan
          \reg{CS}:\reg{IP}.
    \item Nilainya bertambah sendiri setiap kali satu instruksi
          selesai dijalankan.
    \item Programmer hampir tidak pernah menulis \reg{IP} langsung;
          perubahannya terjadi melalui \asm{JMP}, \asm{CALL},
          \asm{RET}, dan \asm{INT}.
\end{itemize}
\end{exampleblock}

\vspace{0.35em}
\begin{alertblock}{Register yang tidak boleh diubah sembarangan}
\small \reg{CS} dan \reg{IP} menentukan instruksi berikutnya;
\reg{SS} dan \reg{SP} menentukan letak \eng{stack}. Bila salah
satunya diubah tanpa alasan benar, prosesor dapat mengambil
instruksi dari alamat salah, menjalankan data sebagai instruksi,
atau membuat program berhenti. \reg{DS} boleh diubah, tetapi harus
dikembalikan sebelum variabel lama diakses.
\end{alertblock}
\end{frame}

% ---------------------------------------------------------
\begin{frame}{Rangkuman Bagian 2}
\small
\begin{itemize}
    \item Register \eng{pointer} dan \eng{index} menyimpan
          \eng{offset}, yaitu jarak sebuah lokasi dari awal segmen.

    \vspace{0.25em}
    \item \reg{SP} menunjuk puncak \eng{stack} dan berubah otomatis
          sebesar 2 pada \asm{PUSH} dan \asm{POP}; \reg{BP} menjadi
          acuan tetap untuk parameter dan variabel lokal.

    \vspace{0.25em}
    \item \reg{SI} menunjuk sumber (\reg{DS}:\reg{SI}) dan \reg{DI}
          menunjuk tujuan (\reg{ES}:\reg{DI}) pada operasi string.

    \vspace{0.25em}
    \item \reg{IP} menyimpan \eng{offset} instruksi berikutnya dan
          selalu berpasangan dengan \reg{CS}.

    \vspace{0.25em}
    \item \reg{CS}, \reg{IP}, \reg{SS}, dan \reg{SP} hanya diubah
          melalui instruksi yang memang dirancang untuk itu.
\end{itemize}
\end{frame}

% =========================================================
% BAGIAN 3
% =========================================================
\PemisahBagian{Register Flag}

% ---------------------------------------------------------
\begin{frame}{Peta Bit Register \reg{FLAGS}}
\small Register \reg{FLAGS} 16 bit --- disebut juga \eng{Program
Status Word} (PSW) --- tidak menyimpan angka hasil hitungan,
melainkan keterangan keadaan hasil operasi terakhir.

\vspace{0.5em}
\begin{center}
\begin{tikzpicture}[
    bit/.style={
        draw=obeBlue, line width=0.7pt, minimum width=0.62cm,
        minimum height=0.60cm, align=center, font=\scriptsize\ttfamily,
        inner sep=0pt
    },
    aktif/.style={bit, fill=obeBlueLight},
    cadangan/.style={bit, fill=obeGrayLight, draw=obeGray},
    nomor/.style={font=\tiny\ttfamily, text=obeGray},
    tanda/.style={font=\tiny\itshape, text=obeGray}
]
    \foreach \i/\nama/\tipe in {
        15/--/cadangan, 14/--/cadangan, 13/--/cadangan,
        12/--/cadangan, 11/OF/aktif, 10/DF/aktif, 9/IF/aktif,
        8/TF/aktif, 7/SF/aktif, 6/ZF/aktif, 5/--/cadangan,
        4/AF/aktif, 3/--/cadangan, 2/PF/aktif, 1/--/cadangan,
        0/CF/aktif}
    {
        \node[\tipe] (b\i) at ({(15-\i)*0.70}, 0) {\nama};
        \node[nomor] at ({(15-\i)*0.70}, -0.48) {\i};
    }
    \node[tanda] at (0, 0.62) {bit tinggi};
    \draw[obeGray, line width=0.6pt, -{Stealth[length=1.6mm]}]
        (0, 0.48) -- (0, 0.34);
    \node[tanda] at (10.5, 0.62) {bit rendah (LSB)};
    \draw[obeGray, line width=0.6pt, -{Stealth[length=1.6mm]}]
        (10.5, 0.48) -- (10.5, 0.34);
    \node[tanda] at (5.25, -1.05)
        {biru = bit dipakai \quad abu-abu = dicadangkan};
\end{tikzpicture}
\end{center}

\vspace{0.2em}
{\footnotesize\color{obeGray}Dari 16 bit yang tersedia, hanya
sembilan bit yang dipakai; sisanya dicadangkan dan nilainya tidak
ditentukan.\par}
\Sumber{Bab V, Peta Bit Register Flag}
\end{frame}

% ---------------------------------------------------------
\begin{frame}{Sembilan Flag: Status dan Kendali}
\footnotesize
\begin{center}
\begin{tabular}{@{}lllp{7.6cm}@{}}
    \toprule
    \textbf{Bit} & \textbf{Flag} & \textbf{Kelompok} &
    \textbf{Arti singkat} \\
    \midrule
    0 & \flag{CF} & Status &
    \eng{Carry}: simpanan keluar dari bit paling tinggi \\
    \addlinespace
    2 & \flag{PF} & Status & \eng{Parity}: paritas byte rendah
    hasil operasi \\
    \addlinespace
    4 & \flag{AF} & Status & \eng{Auxiliary Carry}: simpanan dari
    nibble rendah (bit ke-3) \\
    \addlinespace
    6 & \flag{ZF} & Status & \eng{Zero}: hasil operasi bernilai nol \\
    \addlinespace
    7 & \flag{SF} & Status & \eng{Sign}: bit tanda hasil menyala \\
    \addlinespace
    8 & \flag{TF} & Kendali & \eng{Trap}: satu instruksi per langkah
    untuk \eng{debugging} \\
    \addlinespace
    9 & \flag{IF} & Kendali & \eng{Interrupt}: izin menerima
    \eng{interrupt} tertutup \\
    \addlinespace
    10 & \flag{DF} & Kendali & \eng{Direction}: arah gerak
    \reg{SI}/\reg{DI} pada operasi string \\
    \addlinespace
    11 & \flag{OF} & Status & \eng{Overflow}: hasil bertanda di luar
    jangkauan \\
    \addlinespace
    1, 3, 5, 12--15 & --- & Cadangan & Tidak dipakai \\
    \bottomrule
\end{tabular}
\end{center}
\end{frame}

% ---------------------------------------------------------
\begin{frame}{Enam Flag Status: Kapan Bernilai 1}
\footnotesize
\begin{center}
\begin{tabular}{@{}lp{6.6cm}p{4.6cm}@{}}
    \toprule
    \textbf{Flag} & \textbf{Bernilai 1 bila \dots} &
    \textbf{Contoh instruksi} \\
    \midrule
    \flag{CF} & Operasi \emph{tak bertanda} menghasilkan simpanan
    keluar dari bit tertinggi &
    \asm{ADD}, \asm{SUB}, \asm{CMP}, \asm{ADC}, \asm{SBB},
    \asm{NEG} \\
    \addlinespace
    \flag{PF} & Byte rendah hasil memuat jumlah bit 1 yang
    \emph{genap} & Semua operasi aritmatika dan logika \\
    \addlinespace
    \flag{AF} & Simpanan keluar dari bit ke-3 (nibble rendah) &
    \asm{ADD}, \asm{ADC}, \asm{SUB}, \asm{SBB}, \asm{CMP},
    \asm{INC}, \asm{DEC} \\
    \addlinespace
    \flag{ZF} & Hasil operasi tepat sama dengan nol &
    \asm{ADD}, \asm{SUB}, \asm{CMP}, \asm{AND}, \asm{OR},
    \asm{XOR}, \asm{TEST} \\
    \addlinespace
    \flag{SF} & Bit paling tinggi hasil operasi bernilai 1 &
    \asm{ADD}, \asm{SUB}, \asm{CMP}, \asm{AND}, \asm{OR},
    \asm{XOR}, \asm{NEG} \\
    \addlinespace
    \flag{OF} & Hasil \emph{bertanda} di luar jangkauan:
    $-128\ldots127$ untuk byte & \asm{ADD}, \asm{SUB}, \asm{CMP},
    \asm{IMUL}, \asm{IDIV}, \asm{NEG} \\
    \bottomrule
\end{tabular}
\end{center}

\vspace{0.25em}
{\color{obeGray}\flag{TF}, \flag{IF}, dan \flag{DF} adalah
\eng{flag} kendali: nilainya ditentukan programmer melalui
instruksi seperti \asm{STI}, \asm{CLI}, \asm{STD}, dan
\asm{CLD}.\par}
\Sumber{Bab V, Arti dan Pengaruh Setiap Flag}
\end{frame}

% ---------------------------------------------------------
\begin{frame}{Tiga Kasus \asm{ADD AL, BL}}
\small
\begin{center}
\begin{tabular}{@{}lllcccp{5.5cm}@{}}
    \toprule
    \reg{AL} & \reg{BL} & \reg{AL} baru & \flag{CF} & \flag{OF} &
    \flag{SF} & \textbf{Alasan singkat} \\
    \midrule
    \hex{7F} & \hex{01} & \hex{80} & -- & 1 & 1 &
    Tanpa simpanan keluar, tetapi penjumlahan bertanda melampaui
    batas 127 \\
    \addlinespace
    \hex{80} & \hex{80} & \hex{00} & 1 & 1 & -- &
    Simpanan keluar terjadi; dua bilangan negatif menghasilkan nol \\
    \addlinespace
    \hex{FF} & \hex{01} & \hex{00} & 1 & -- & -- &
    Simpanan keluar terjadi, tetapi hasil bertanda tepat nol \\
    \bottomrule
\end{tabular}
\end{center}

\vspace{0.35em}
{\footnotesize
\begin{itemize}
    \item \hex{7F} $+$ \hex{01}: dua operand positif menghasilkan
          hasil yang ditafsirkan negatif $\rightarrow$ \flag{OF} $=1$.
    \item \hex{80} $+$ \hex{80}: \dec{128} $+$ \dec{128} $=$
          \dec{256} $\rightarrow$ \flag{CF} $=1$ dan \flag{ZF} $=1$.
    \item \hex{FF} $+$ \hex{01}: \hex{FF} adalah $-1$ dan \hex{01}
          adalah $+1$, sehingga hasil bertandanya nol $\rightarrow$
          \flag{OF} $=0$.
\end{itemize}\par}
\Sumber{Bab V, Contoh Perubahan Flag pada ADD AL, BL}
\end{frame}

% ---------------------------------------------------------
\begin{frame}{\flag{CF} dan \flag{OF} Menjawab Pertanyaan Berbeda}
\begin{alertblock}{Jangan sampai tertukar}
\small \flag{CF} berlaku untuk bilangan \textbf{tak bertanda},
sedangkan \flag{OF} berlaku untuk bilangan \textbf{bertanda}.
Memakai \asm{JC} untuk memeriksa hasil bilangan bertanda adalah
salah; yang tepat \asm{JO} atau \asm{JS}.
\end{alertblock}

\vspace{0.4em}
\begin{block}{\asm{MOV} dan kawan-kawan tidak menyentuh \eng{flag}}
\small \asm{MOV}, \asm{LEA}, \asm{PUSH}, dan \asm{POP} \textbf{tidak}
mengubah \eng{flag}. Pengecualiannya \asm{POPF} dan \asm{SAHF} yang
memang dirancang untuk menulis \eng{flag}.
\end{block}

\vspace{0.4em}
\begin{exampleblock}{\asm{INC} dan \asm{DEC} juga bukan pengubah penuh}
\small Keduanya mengubah \flag{ZF}, \flag{SF}, \flag{OF}, dan
\flag{AF}, tetapi membiarkan \flag{CF} apa adanya. Untuk menambah 1
sekaligus memperbarui \flag{CF}, pakai \code{ADD AL, 1}.
\end{exampleblock}
\end{frame}

% ---------------------------------------------------------
\begin{frame}{Flag Dibaca oleh Lompatan Bersyarat}
\footnotesize \eng{Flag} baru bermakna ketika dibaca lompatan
bersyarat. \asm{CMP} membandingkan dua nilai tanpa menyimpan hasil,
lalu meninggalkan jejaknya pada \eng{flag}.

\vspace{0.3em}
\begin{center}
\begin{tabular}{@{}lp{4.4cm}p{6.6cm}@{}}
    \toprule
    \textbf{Flag} & \textbf{Lompatan} & \textbf{Makna} \\
    \midrule
    \flag{ZF} & \asm{JE}, \asm{JZ} &
    Lompat bila hasil nol (kedua nilai sama) \\
    \addlinespace
    \flag{ZF} & \asm{JNE}, \asm{JNZ} &
    Lompat bila hasil tidak nol (kedua nilai berbeda) \\
    \addlinespace
    \flag{CF} & \asm{JC}, \asm{JB} &
    Lompat bila ada simpanan keluar (nilai pertama lebih kecil secara
    tak bertanda) \\
    \addlinespace
    \flag{CF} & \asm{JNC}, \asm{JAE} &
    Lompat bila tidak ada simpanan keluar \\
    \addlinespace
    \flag{SF} & \asm{JS}, \asm{JNS} &
    Lompat bila bit tanda menyala atau tidak \\
    \addlinespace
    \flag{OF} & \asm{JO}, \asm{JNO} &
    Lompat bila terjadi \eng{overflow} bertanda atau tidak \\
    \addlinespace
    \flag{PF} & \asm{JP}, \asm{JNP} &
    Lompat berdasarkan paritas byte rendah hasil \\
    \bottomrule
\end{tabular}
\end{center}
\Sumber{Bab V, Pengaruh Flag terhadap Keputusan Program}
\end{frame}

% ---------------------------------------------------------
\begin{frame}{Rangkuman Bagian 3}
\small
\begin{itemize}
    \item \reg{FLAGS} 16 bit memuat sembilan \eng{flag} aktif: enam
          \eng{flag} status (\flag{CF}, \flag{PF}, \flag{AF},
          \flag{ZF}, \flag{SF}, \flag{OF}) dan tiga \eng{flag}
          kendali (\flag{TF}, \flag{IF}, \flag{DF}).

    \vspace{0.25em}
    \item \flag{CF} melaporkan simpanan keluar bilangan tak
          bertanda; \flag{OF} melaporkan \eng{overflow} bilangan
          bertanda. Keduanya tidak boleh dipertukarkan.

    \vspace{0.25em}
    \item \flag{ZF} menandai hasil nol, \flag{SF} bit tanda yang
          menyala, \flag{PF} paritas byte rendah, dan \flag{AF}
          simpanan keluar dari nibble rendah.

    \vspace{0.25em}
    \item \asm{MOV}, \asm{LEA}, \asm{PUSH}, dan \asm{POP} tidak
          mengubah \eng{flag}; \asm{INC} dan \asm{DEC} mengubah
          semuanya kecuali \flag{CF}.

    \vspace{0.25em}
    \item \eng{Flag} dipakai lompatan bersyarat: \asm{JE}/\asm{JNE},
          \asm{JC}/\asm{JNC}, \asm{JS}/\asm{JNS}, \asm{JO}/\asm{JNO},
          dan \asm{JP}/\asm{JNP}.
\end{itemize}
\end{frame}

% =========================================================
% CONTOH PROGRAM
% =========================================================
\begin{frame}[fragile]{Contoh 1: Memprediksi Flag Sebelum
Menjalankan Program}
\small Dua penjumlahan yang sengaja dipilih untuk memicu kondisi
\eng{flag} berbeda.
\vspace{0.2em}
\begin{columns}[t]
\begin{column}{0.60\textwidth}
\begin{kodeasm}[numbers=none]
.DATA
    hasil1 DW 0
    hasil2 DW 0
.CODE
START:
    MOV AX, @DATA
    MOV DS, AX
    MOV AX, 007Fh      ; kasus 1: +127
    MOV BX, 0001h
    ADD AX, BX         ; AX=0080h: SF=1, OF=1
    MOV hasil1, AX     ; MOV tak mengubah flag
    MOV AX, 8000h      ; kasus 2: -32768
    MOV BX, 8000h
    ADD AX, BX         ; AX=0000h: CF=1, ZF=1
    MOV hasil2, AX
\end{kodeasm}
\end{column}
\begin{column}{0.36\textwidth}
{\color{obeTeal}\bfseries Keluaran}\\[0.3em]
\begin{keluaran}
Pengamatan flag
selesai.
\end{keluaran}

\vspace{0.4em}
{\footnotesize\flag{CF}, \flag{OF}, \flag{SF}, dan \flag{ZF}
diprediksi dahulu, lalu diperiksa dengan \eng{Single Step}.\par}

\vspace{0.3em}
{\footnotesize\color{obeGray}\file{prediksi-flag.asm}\par}
\end{column}
\end{columns}
\end{frame}

% ---------------------------------------------------------
\begin{frame}[fragile]{Contoh 2: Memilih Pesan Berdasarkan Flag}
\small \hex{F0} ditambah \hex{10} menghasilkan \hex{100}: \reg{AL}
menyimpan \hex{00} dengan \flag{CF} $=1$ dan \flag{ZF} $=1$.
\vspace{0.2em}
\begin{columns}[t]
\begin{column}{0.58\textwidth}
\begin{kodeasm}[numbers=none]
    MOV AL, 0F0h       ; AL = 1111 0000b
    MOV BL, 10h        ; BL = 0001 0000b
    ADD AL, BL         ; AL=00h, CF=1, ZF=1

    JNC tanpa_carry    ; lompat bila CF = 0
    MOV AH, 09h        ; cetak pesan carry
    LEA DX, pesan_carry
    INT 21h
tanpa_carry:
    JNZ bukan_nol      ; lompat bila ZF = 0
    MOV AH, 09h        ; cetak pesan hasil nol
    LEA DX, pesan_nol
    INT 21h
bukan_nol:
    JNS tanpa_negatif  ; lompat bila SF = 0
\end{kodeasm}
\end{column}
\begin{column}{0.38\textwidth}
{\color{obeTeal}\bfseries Keluaran}\\[0.3em]
\begin{keluaran}
Terjadi carry,
CF = 1.
Hasil penjumlahan
nol, ZF = 1.
\end{keluaran}

\vspace{0.4em}
{\footnotesize Layanan \code{INT 21h} tidak merusak \eng{flag}:
\asm{INT} menyimpan \reg{FLAGS} ke \eng{stack} dan \asm{IRET}
mengembalikannya.\par}

\vspace{0.3em}
{\footnotesize\color{obeGray}\file{pesan-flag.asm}\par}
\end{column}
\end{columns}
\end{frame}

% ---------------------------------------------------------
\begin{frame}[fragile]{Contoh 3: Akses Byte Rendah dan Byte Tinggi}
\small Satu register 16 bit dapat diperlakukan sebagai dua register
8 bit yang mandiri.
\vspace{0.2em}
\begin{columns}[t]
\begin{column}{0.58\textwidth}
\begin{kodeasm}[numbers=none]
.MODEL SMALL
.STACK 100h

.DATA
    dummy DB 0          ; agar segmen data ada

.CODE
START:
    MOV AX, @DATA
    MOV DS, AX

    MOV AX, 1234h       ; AH = 12h, AL = 34h
    MOV BL, AH          ; BL = 12h
    MOV BH, AL          ; BH = 34h
    ADD AL, BL          ; AL = 46h, AH tetap
    MOV CL, AL          ; CL = 46h
\end{kodeasm}
\end{column}
\begin{column}{0.38\textwidth}
{\color{obeTeal}\bfseries Keadaan akhir}\\[0.3em]
\begin{keluaran}
AX = 1246h
BX = 3412h
CL = 46h
\end{keluaran}

\vspace{0.5em}
{\footnotesize \code{ADD AL, BL} hanya mengubah \reg{AL}. Nilai
\reg{AH} tetap \hex{12} karena operasi 8 bit pada \reg{AL} tidak
menyentuh byte tinggi.\par}
\end{column}
\end{columns}
\end{frame}

% =========================================================
% PRAKTIKUM
% =========================================================
\begin{frame}{Praktikum: Pengamatan Register dan Flag}
\begin{block}{Judul praktikum}
\small Pengamatan register dan \eng{flag} dengan \eng{Single Step}
pada Emu8086 \quad\textbullet\quad \textbf{Sub-CPMK:} 1.2
\quad\textbullet\quad \textbf{Waktu:} 100 menit
\quad\textbullet\quad \textbf{Bentuk:} individu
\end{block}

\vspace{0.3em}
\begin{columns}[t]
\begin{column}{0.48\textwidth}
{\color{obeBlue}\bfseries Alat dan bahan}
\begin{itemize}\small
    \item Komputer dengan Emu8086 terpasang
    \item Berkas \file{prediksi-flag.asm} dan
          \file{pesan-flag.asm}
    \item Lembar kerja pengamatan register dan \eng{flag}
\end{itemize}
\end{column}
\begin{column}{0.48\textwidth}
{\color{obeBlue}\bfseries Langkah kerja}
\begin{enumerate}\small
    \item Buka jendela \file{Registers} dan \file{Flags}, lalu
          jalankan dengan \eng{Single Step}
    \item Catat \reg{AX}, \reg{BX}, \reg{SP}, dan \reg{IP}; amati
          \eng{flag} sebelum dan sesudah \code{ADD AX, BX}
    \item Ubah \reg{AX} menjadi \hex{FFFF} dan \reg{AL} menjadi
          \hex{00} secara manual
    \item Ubah \flag{CF} dan \flag{ZF} pada
          \file{pesan-flag.asm}, amati pesan yang tercetak
\end{enumerate}
\end{column}
\end{columns}

\vspace{0.3em}
{\footnotesize\textbf{Dikumpulkan:} lembar pengamatan register,
lembar prediksi \eng{flag}, tangkapan layar \file{Registers} dan
\file{Flags}, serta catatan percobaan mengubah \flag{CF} dan
\flag{ZF}.\par}
\end{frame}

% =========================================================
% LATIHAN
% =========================================================
\begin{frame}{Latihan}
\small
\begin{enumerate}
    \item Jelaskan apa yang dimaksud dengan register, lalu sebutkan
          empat kelompok register 8086 beserta satu contoh register
          pada setiap kelompok.

    \vspace{0.3em}
    \item \reg{AX} berisi \hex{ABCD}. Tentukan nilai \reg{AH} dan
          \reg{AL} setelah \code{MOV AL, 55h} dijalankan. Mengapa
          \reg{AH} tidak berubah?

    \vspace{0.3em}
    \item Seorang programmer ingin menjumlahkan dua nilai,
          menyimpan hasilnya, lalu mengulangi proses sepuluh kali.
          Register mana yang tepat untuk menampung hasil dan mana
          yang tepat menjadi pencacah? Jelaskan alasannya.

    \vspace{0.3em}
    \item Untuk \code{ADD AL, BL}, tentukan \flag{CF}, \flag{OF},
          \flag{SF}, dan \flag{ZF} bila \reg{AL} dan \reg{BL}
          bernilai (\emph{a}) \hex{7F} dan \hex{01},
          (\emph{b}) \hex{05} dan \hex{05}. Sertakan penjelasan
          singkat.

    \vspace{0.3em}
    \item Sebuah program berisi \code{MOV SP, 0000h} sebelum
          memanggil prosedur. Jelaskan mengapa berbahaya dan apa
          yang terjadi saat \asm{RET} dijalankan.
\end{enumerate}
\end{frame}

% =========================================================
% RANGKUMAN
% =========================================================
\begin{frame}{Rangkuman}
\small
\begin{itemize}
    \item Register adalah tempat penyimpanan tercepat di dalam
          prosesor. Prosesor 8086 memiliki empat belas register
          16 bit yang dapat dipakai programmer.

    \vspace{0.25em}
    \item \reg{AX}, \reg{BX}, \reg{CX}, dan \reg{DX} dapat diakses
          sebagai dua byte terpisah. Menulis ke byte rendah tidak
          mengubah byte tinggi.

    \vspace{0.25em}
    \item Peran khas muncul dari cara instruksi disandikan:
          \reg{AX} akumulator, \reg{BX} basis pengalamatan,
          \reg{CX} pencacah, \reg{DX} pasangan data dan penunjuk
          port.

    \vspace{0.25em}
    \item Register segmen menyimpan nomor segmen, bukan data;
          pasangannya \reg{IP} untuk kode, \reg{SP}/\reg{BP} untuk
          \eng{stack}, serta \reg{SI} dan \reg{DI} untuk data.

    \vspace{0.25em}
    \item \reg{FLAGS} memuat sembilan \eng{flag} aktif: enam status
          dan tiga kendali. \flag{CF} untuk bilangan tak bertanda,
          \flag{OF} untuk bilangan bertanda.

    \vspace{0.25em}
    \item \reg{CS}, \reg{IP}, \reg{SS}, dan \reg{SP} mengendalikan
          alur program dan letak \eng{stack}, sehingga tidak boleh
          diubah sembarangan.
\end{itemize}
\end{frame}

% ---------------------------------------------------------
\begin{frame}{Cek Pemahaman}
\small
\begin{enumerate}
    \item Setelah \code{MOV AL, 20h} dijalankan pada \reg{AX} yang
          semula berisi \hex{1234}, isi \reg{AX} menjadi \dots
          \begin{enumerate}\footnotesize
              \item \hex{1220} \quad \item \hex{2034} \quad
              \item \hex{0020} \quad \item \hex{2020}
          \end{enumerate}

    \vspace{0.3em}
    \item Register yang selalu berpasangan dengan \reg{CS} untuk
          menentukan alamat instruksi berikutnya adalah \dots
          \begin{enumerate}\footnotesize
              \item \reg{SP} \quad \item \reg{IP} \quad
              \item \reg{DS} \quad \item \reg{DI}
          \end{enumerate}

    \vspace{0.3em}
    \item \code{ADD AL, BL} dengan \reg{AL} $=$ \hex{7F} dan
          \reg{BL} $=$ \hex{01} menghasilkan \dots
          \begin{enumerate}\footnotesize
              \item \flag{CF} $=1$ dan \flag{OF} $=0$
              \item \flag{CF} $=0$ dan \flag{OF} $=1$
              \item \flag{CF} $=1$ dan \flag{OF} $=1$
              \item \flag{CF} $=0$ dan \flag{OF} $=0$
          \end{enumerate}
\end{enumerate}
\end{frame}

% =========================================================
% PENUTUP
% =========================================================
\begin{frame}{Penutup}
\small
\begin{block}{Sebelum pertemuan berikutnya}
\begin{itemize}
    \item Selesaikan praktikum: lembar pengamatan register dan
          lembar prediksi \eng{flag}.
    \item Tuliskan potongan program yang memuat \hex{F0} ke
          \reg{AL} dan \hex{10} ke \reg{BL}, lalu cetak ``NOL'' bila
          \flag{ZF} bernilai 1.
    \item Baca \textbf{Bab VI} --- Segmentasi Memori dan Alamat
          Fisik.
\end{itemize}
\end{block}

\vspace{0.5em}
\begin{center}
{\color{obeBlue}\sffamily\large\bfseries Terima kasih\par}
\vspace{0.2em}
{\color{obeGray}\sffamily\small \NamaMK\ (\KodeMK) \textbullet\
\SemesterTeks\par}
\end{center}
\end{frame}

\end{document}
